\subsection*{An exact common-kernel counterexample}

\statement{B.57}{Refuted}{P7: the three resolvent kernels have a strictly larger intersection}
With the corrected third target of B.53, the proposed equality
\[
 \mathcal K=
 \ker\mathfrak T_+\cap\ker\mathfrak T_-\cap\ker\mathfrak T_2
                         \quad\hbox{inside }\mathcal H
\]
is false. There is a nonempty specified set of full peripheral
lifts $\mathscr C\subset\mathcal H$ such that every
$\Phi\in\mathscr C$ satisfies
\[
 \mathfrak T_+(\Phi)=\mathfrak T_-(\Phi)
      =\mathfrak T_2(\Phi)=1,\qquad \Phi\notin\mathcal K.     \tag{B.57.1}
\]
Each such element has a nontrivial tripod action, detected by
the outer automorphism of $A_9$ induced by an odd transposition.
In particular the common kernel is infinite and contains a
closed copy of $\mathbb Z_2$. This decides its nontriviality,
but does not decide whether $\mathcal K$ is trivial.

\begin{proof}
We use an infinite retained quotient to impose all three
resolvent identities exactly. A finite alternating-group
quotient supplies the nongeometric witness. The central
refinement is constructed explicitly before applying the
closed-normal lifting criterion of A.3 and A.34.

\paragraph{The retained quotient and the fifteen normal kernels.}
Use $J=\widehat{\Gamma_T}=C_2\amalg C_3\amalg C_3$ from B.53
and denote its three filling maps by
\[
 \rho_+,\rho_-:J\twoheadrightarrow\widehat B,\qquad
 \rho_2:J\twoheadrightarrow D_2=C_3\amalg C_3.
\]
Let $L_i=\widehat K\cap\ker\rho_i$, for
$i\in\{+,-,2\}$. These are closed normal subgroups of
$\widehat K$, invariant under conjugation by $J$.
Take one representative of each of the five cosets of $J$
in $\widehat B$ and put
\[
 V=\bigcap_{g\in\widehat B/J}\ \bigcap_{i\in\{+,-,2\}}
                              gL_i g^{-1},\qquad
 C=\widehat B/V.                                          \tag{B.57.2}
\]
Conjugation permutes the fifteen subgroups in this
intersection. Thus $V$ is closed and normal in $\widehat B$,
and $V\leq\widehat K$. In particular $C$ still maps onto
$A_5$. Any automorphism of $\widehat B$ inducing the identity
on $C$ induces the identity on all three filled groups:
for $j\in J$, its image differs from $j$ by an element of
$V\subseteq L_i$. We shall retain this stronger, based
identity, not just finitely many target tests.

Here are the exact images and cusp distributions for these
fillings. In the finite quotient $J/\widehat K=A_4$,
the normal closure of either order-three factor is all
of $A_4$. Therefore
\[
                 \rho_+(\widehat K)=\widehat B
                      =\rho_-(\widehat K).                \tag{B.57.3}
\]
The cover $\widehat K\leq J$ has degree twelve, with twelve
cusps of width one relative to the unique cusp of $J$.
In either image (B.57.3), their meridians are consequently
conjugates of the one target cusp meridian.

The normal closure of the order-two factor in $A_4$ is
its $V_4$. Hence
\[
 \rho_2(\widehat K)=P_2
   =\ker\bigl(C_3\amalg C_3\longrightarrow C_3,\
                        a\mapsto1,\ b\mapsto-1\bigr).
                                                               \tag{B.57.4}
\]
This is the full free three-cusp group of rank two.
The twelve cusps of $\widehat K$ map to its three cusp
classes in blocks of four, each meridian with exponent
one. One may see the blocks either from the $V_4$ fibres
of $A_4\to C_3$, or from the index-three covering of the
$(3,3,\infty)$ orbifold. Both (B.57.3) and (B.57.4)
also hold for all five conjugates of the cover.

\paragraph{A specified alternating-group quotient of a tripod.}
Choose a three-subset $t$ of the twelve vertices consisting
of an antipodal pair and one further vertex. Its stabilizer
in $A_5$ is trivial. Indeed a permutation preserving $t$
must preserve its unique antipodal pair; the normalizer
$D_5$ of the corresponding $C_5$ acts freely on the other
ten vertices. The orbit of $t$ therefore has sixty members.
This also follows directly from the twelve-coset
permutation certificate described below.

Let $\pi_t:\widehat K\twoheadrightarrow F_t$ kill the
other nine cusp classes and choose positive meridians
$x,y,z=(xy)^{-1}$ for the resulting tripod. In the
natural nine-letter action, with the rightmost permutation
applied first, put
\[
 X=(1\,2\,3\,4\,5\,6\,7),\qquad
 Y=(1\,2\,5\,4\,3\,8\,9),\qquad Z=(XY)^{-1}.
                                                               \tag{B.57.5}
\]
The exact enumeration in
\texttt{work\_CE/code/check\_common\_kernel\_alt9.py}
proves
\[
 \langle X,Y\rangle=A_9,\qquad
        X,Y,Z\hbox{ all have cycle type }(7,1,1).           \tag{B.57.6}
\]
Thus the map
\[
       q_0:\widehat K\xrightarrow{\pi_t}F_t
                     \twoheadrightarrow S=A_9,
                  \qquad x\mapsto X,\quad y\mapsto Y      \tag{B.57.7}
\]
is onto. It kills precisely nine of the twelve cusp
classes, and is nontrivial on each of the three retained
ones.

Every one of the fifteen normal subgroups in (B.57.2)
has image $S$ under $q_0$. Since $S$ is simple, a proper
image would be trivial, so $q_0$ would factor through
the corresponding filled image. For a filling of an
order-three factor, (B.57.3) would then identify all
twelve cusp images up to conjugacy. A killed meridian
would kill the target cusp $\tau$; since
$\widehat B/\overline{\langle\!\langle\tau\rangle\!\rangle}=1$,
the entire factor map would be trivial, contrary to
(B.57.6). For the order-two filling, every one of the
three blocks of four in (B.57.4) contains a killed
meridian: only three meridians are retained in total.
Factoring would therefore kill all three peripheral
generators of $P_2$, again a contradiction.
The argument is unchanged by any of the five
conjugations.

Write these fifteen normal subgroups as $N_1,\ldots,N_{15}$.
The iterated closed commutator
\[
        [\,\cdots[[N_1,N_2],N_3],\ldots,N_{15}\,]
\]
is contained in their intersection $V$ and has image
$S$ under $q_0$, because $[S,S]=S$ at every step.
Consequently
\[
                              q_0(V)=S.                  \tag{B.57.8}
\]
This is a direct commutator argument. It does not use
an invalid interchange of image and intersection.

\paragraph{The induced quotient has its entire perfect base.}
Induce (B.57.7) from the normal subgroup $\widehat K$
of index sixty. The usual coset formula gives a
continuous homomorphism
\[
 q:\widehat B\longrightarrow S^{A_5}\rtimes A_5.
\]
For example, with representatives $t_\omega$ mapping
to $\omega\in A_5$, its base coordinates are
$q_0(t_\omega b\,t_{\omega\bar b}^{-1})$, and the
permutation coordinate is $\bar b=q_5(b)$; this formula
uses the regular right action on the coordinates, and
we choose $t_1=1$.
On $\widehat K$, the sixty coordinate kernels are
distinct. Their nonzero cusp supports are the sixty
distinct triples in the orbit of $t$.

For clarity, distinct kernels suffice for the complete
product here. Given one coordinate $j$, each of the
other coordinate kernels has image $S$ in coordinate
$j$; otherwise that maximal normal subgroup would
equal the $j$th kernel. Their iterated commutator is
contained in all the other kernels and still maps
onto $S$ in coordinate $j$. Thus the image contains
the whole $j$th factor alone. Doing this for every
$j$ proves
\[
 q(\widehat K)=M=S^{60},\qquad
 q(\widehat B)=E=M\rtimes A_5.                             \tag{B.57.9}
\]
The latter equality follows because the image contains
the entire base and surjects onto the permutation group.
This is the full regular wreath product, not an
unspecified subdirect subgroup.

By (B.57.8), $q(V)$ projects onto the coordinate
corresponding to $q_0$ and, by normality and transitivity,
onto all sixty coordinates. It is normal in $E$ and
therefore normal in $M$. Normal subgroups of a product
of nonabelian finite simple groups are products of
factors. It follows that
\[
                               q(V)=M.                   \tag{B.57.10}
\]
One can also prove the last conclusion by commuting
with each individual factor of $M$: a surjective
projection of the normal subgroup supplies that
entire factor.

Put $N=V\cap\ker q$. Equation (B.57.10) proves that
the joint quotient is exactly
\[
 G=\widehat B/N
       =C\times_{A_5}E
       \cong M\rtimes C,                                  \tag{B.57.11}
\]
where $C$ acts through $A_5$ by regular coordinate
permutations. To check surjectivity directly, lift
$c\in C$ to $b\in\widehat B$ and correct the difference
between $q(b)$ and the specified compatible element
of $E$ by an element of $V$, using (B.57.10).
The section $C\to G$ in (B.57.11) is
$c\mapsto(c,(1,\bar c))$.

\paragraph{The finite automorphism and its lifted classes.}
Let $\eta_0$ be conjugation on $S=A_9$ by the
transposition $(1\,2)\in S_9$. It is outer: an inner
realization would make an odd permutation centralize
$A_9$, whereas that centralizer in $S_9$ is trivial.
The certificate checks this last assertion already
on the specified generators $X,Y$.
Apply $\eta_0$ to every factor of $M$ and the identity
to $A_5$. This gives an automorphism $\eta$ of $E$,
and, by (B.57.11), an automorphism
\[
                     \alpha=\eta_M\rtimes\operatorname{id}_C
                                     \quad\hbox{of }G.     \tag{B.57.12}
\]

The one finite-group input about a universal cover is
\[
                H_2(A_9,\mathbb Z)=C_2,\qquad
                         \widetilde S=2.A_9.
                                                               \tag{B.57.13}
\]
This is the Schur multiplier and double cover recorded
in the primary ATLAS entry for $A_9$.\cite{CEAtlasA9}
An automorphism lifts uniquely to the universal central
cover of a perfect group. Hence $\eta_0$ lifts to
$\widetilde\eta_0$ on $\widetilde S$ and fixes its central
$C_2$ pointwise. The universal central cover of $M$ is
\[
                       \widetilde M=(2.A_9)^{60}.
\]
For a product of perfect groups, preimages of different
factors commute in every central extension: their
commutator is central and factors through the
abelianizations, which vanish. This proves the
displayed product assertion directly. The diagonal
lift $\widetilde\eta_M$ consequently fixes the full
central kernel $C_2^{60}$ pointwise.

Every seven-cycle in $A_9$ is conjugate in $A_9$ to
its $\eta_0$-image. Indeed it has two fixed letters;
their transposition is an odd element of its
$S_9$-centralizer. Multiplying $(1\,2)$ by that
centralizer element produces an even conjugator.
Moreover $\widetilde\eta_0$ preserves the conjugacy
class of \emph{each} element over a seven-cycle.
Of the two lifts of that cycle, exactly one has
order seven. Conjugacy of the base elements and
uniqueness of this odd-order lift give conjugacy
of the order-seven lifts. Multiplication by the
fixed central involution gives the same result
for the other lift. The two elements over the
identity are central and fixed as well. Thus
\[
 \widetilde\eta_0\hbox{ preserves every lift class
                 over a seven-cycle or the identity}.   \tag{B.57.14}
\]

\paragraph{A central-extension lemma with no finite-index assumption.}
Consider the actual central refinement
\[
 \mathscr E=\widehat B/[N,\widehat B],\qquad
 Z_N=N/[N,\widehat B],\qquad
                   1\to Z_N\to\mathscr E\to G\to1.
                                                               \tag{B.57.15}
\]
All commutator subgroups here are closed.
Let $\mathscr M$ be the preimage of $M$.
The universal property of $\widetilde M$ gives a
canonical homomorphism $\widetilde M\to\mathscr M$;
let $M^\#$ be its image. It is a finite perfect group
mapping onto $M$, and
\[
                   \mathscr M=M^\# Z_N.
\]
These facts do not require $\mathscr E$ or $C$ to
be finite. For example the universal property follows
from the free-presentation construction
$[F,F]/[F,R]\to M$ of the universal central cover:
lift the free generators into the central extension
and restrict to the commutator subgroup. Independence
and uniqueness follow from perfectness.

Let $C^\#$ be the preimage in $\mathscr E$ of the
specified section $C\leq G$. Then
\[
       \mathscr E=M^\# C^\#,\qquad
       M^\#\cap C^\#\subseteq Z_N.                         \tag{B.57.16}
\]
The kernel of $\widetilde M\to M^\#$ is central.
Since $\widetilde\eta_M$ fixes the entire universal
central kernel, it descends to an automorphism
$\eta^\#$ of $M^\#$, fixing $M^\#\cap C^\#$.
It also commutes with the conjugation action of
$C^\#$ on $M^\#$. On $M$ that action is a factor
permutation, with which the diagonal $\eta_M$
commutes. Its lift to the perfect cover is unique:
two lifts differ by a homomorphism from a perfect
group to a central abelian group, hence coincide.
It follows that
\[
              \alpha^\#(mc)=\eta^\#(m)c
                    \quad(m\in M^\#,\,c\in C^\#)           \tag{B.57.17}
\]
is a well-defined continuous automorphism of
$\mathscr E$ covering (B.57.12). Equation (B.57.16),
the fixed intersection, and the commuting actions
verify respectively its definition, its independence
of the decomposition, and its multiplication law.
Its inverse is obtained in the same way.

\paragraph{The three actual classes in the central refinement.}
It remains to verify the precise conjugacy conditions
in $\mathscr E$, since an automorphism of an abstract
convenient central cover of $E$ would not suffice.

In the regular wreath product (B.57.9), the image of
$s$ is conjugate by $M$ to the pure permutation
$\bar s$, and the image of $r$ is conjugate by $M$
to the pure permutation $\bar r$. Indeed their
permutations have respectively only two-cycles and
three-cycles, and the cycle products in the base
are all one because $s^2=r^3=1$. Lifting these
conjugators to $M^\#$ conjugates the corresponding
elements of $\mathscr E$ into $C^\#$. They are
then fixed by (B.57.17). Thus $\alpha^\#$ sends
both actual elements $s,r$ of $\mathscr E$ to
conjugates by elements of $M^\#$.

Write the actual cusp in $\mathscr E$ as
$\tau=m^\#c^\#$, with $m^\#\in M^\#$ and
$c^\#\in C^\#$. Conjugation by $c^\#$ on $M^\#$
is the lift of the permutation $\bar\tau\in A_5$,
of order five. Choose a lift $\widetilde m$ of
$m^\#$ to $\widetilde M$. Twisted conjugacy in a
permutational wreath product is determined by the
conjugacy classes of the products around each
permutation cycle. This follows by solving the
coordinate conjugator equation successively along
each cycle; its last equation is exactly conjugacy
of the cycle product.

The twelve cycles of $\bar\tau$ on the sixty
coordinates correspond to the twelve cusp classes
of $\widehat K$. Their base cycle products in $M$
are the $q_0$-images of conjugates of $\tau^5$.
By (B.57.7) each is either the identity or a
seven-cycle. The corresponding products of
$\widetilde m$ in $\widetilde M$ therefore lie over
precisely those elements. By (B.57.14), every such
product is conjugate to its $\widetilde\eta_0$-image.
The coordinate equations give
$\widetilde u\in\widetilde M$ such that
\[
 \widetilde\eta_M(\widetilde m)
       =\widetilde u\,\widetilde m\,
                              \bar\tau(\widetilde u^{-1}).
\]
Projecting to $M^\#$, and recalling the action of
$c^\#$, yields
\[
                         \alpha^\#(\tau)=u\tau u^{-1}.
                                                               \tag{B.57.18}
\]
This calculation takes place above the actual source
element $\tau$. It is still valid if the central
quotient defining $M^\#$ couples different factors,
because it was first solved in the full product
$\widetilde M$ and then projected.

\paragraph{The full lift and its nongeometric tripod.}
Equations (B.57.17)--(B.57.18) give an automorphism
of the actual group $\widehat B/[N,\widehat B]$
preserving the conjugacy classes of $s,r,\tau$,
with literal conjugators there. The closed-normal
modular version of the central-quotient criterion
in A.3 and A.34 therefore gives a continuous
automorphism $\Phi$ of $\widehat B$ inducing
$\alpha$ on $\widehat B/N$, with those conjugators
projected to $\widehat B/N$ and then lifted.
The criterion does not assert that the full lift
retains the entire chosen automorphism $\alpha^\#$
of the central refinement; only its quotient
automorphism $\alpha$ and the projected conjugators
are needed here.
Let $\mathscr C$ be this nonempty set of
full lifts, regarded in the peripheral outer group.
By B.6 its elements belong to $\mathcal H$.

Their action on $C=\widehat B/V$ is the identity.
The observation after (B.57.2) gives the three exact
identities in (B.57.1), with the target types and
geometric markings of B.53. The action on $A_5$ is
also the identity, and the multiplier is one, so
the individual cusp classes of $\widehat K$ are
fixed and the tripod quotient $\pi_t$ is preserved.
On the identity coordinate of (B.57.9),
\[
                   q_0\circ\Phi|_{\widehat K}
                                =\eta_0\circ q_0.         \tag{B.57.19}
\]
If the induced outer action on $F_t$ were trivial,
a representative would be conjugation by an element
of $F_t$. Applying its quotient map onto $A_9$
would make $\eta_0$ inner, contradicting its
construction. Thus this tripod action is nontrivial.
B.10 says every element of $\mathcal K$ has trivial
tripod actions, proving $\Phi\notin\mathcal K$.

In particular $\Phi\ne1$. We specify the finite outer
quotient used in the final assertion. Consider the
closed subgroup of peripheral outer classes admitting
a representative which induces the identity on $A_5$
and preserves $\ker q_0\leq\widehat K$.
Two such representatives of the same outer class differ
by conjugation by an element of $\widehat K$: its image
must centralize $A_5$, which has trivial center.
This ambiguity induces an inner automorphism of $A_9$.
There is therefore a well-defined continuous map from
that closed subgroup to $\operatorname{Out}(A_9)$.
It takes $\Phi$ to the nontrivial outer transposition,
and it is defined on the entire closed cyclic subgroup
$L=\overline{\langle\Phi\rangle}$.
R1 makes $\mathcal H$ torsion-free, so $L$ is infinite.
The indicated finite quotient of $L$ has order two.
A torsion-free procyclic profinite
group with a quotient of even order has a closed
factor isomorphic to $\mathbb Z_2$. This factor
lies in the common kernel, which proves the last
assertion of the statement. Such a $\Phi$ is also
outside the Galois image: its multiplier is one,
so the arithmetic identity in B.53 would make
$\mathfrak T_+(\Phi)=\Phi$ there, in conflict with
(B.57.1) and $\Phi\ne1$.
\end{proof}

The finite certificate enumerates all $181440$ elements
generated by (B.57.5), checks the three cycle types and
explicit even conjugators, checks the trivial joint
centralizer in $S_9$, and verifies the sixty-element
icosahedral tripod orbit and the finite $A_4$ normal
closures. The universal multiplier (B.57.13) is the
separately cited primary finite-group input. The full
counterexample is the nonempty lift set defined by
(B.57.2), (B.57.7), (B.57.11), and (B.57.17); no
limiting compatibility is inferred from a finite
sample. The strict containment in (B.57.1) leaves
the original centralizer problem for $\mathcal K$
undecided, and supplies no pentagon assertion.

